\section{Experimental Setup and Results}
\label{sec:exp}

\subsection{Setup}

The current comparison table reports seven benchmarks and five methods. The experiments compare the proposed timing-driven dynamic-programming (DP) flow against Vivado, \rsad, AMFPlace 2.0 + Vivado P\&R, and DSPlace + Vivado P\&R. The target audience should read the first two comparisons differently:
\begin{itemize}[leftmargin=*,topsep=2pt,itemsep=2pt]
\item Vivado is the industrial baseline.
\item \rsad is the direct predecessor and isolates the value of the new timing-driven formulation over the prior OC-based HPWL-oriented deterministic method.
\end{itemize}

This distinction is important. Relative to Vivado, the paper should emphasize that problem-specific systolic-array placement already has large value. Relative to \rsad, the paper should emphasize that the new contribution is not better HPWL minimization alone, but a better timing tradeoff through explicit control of the longest wirelength.

\subsection{Main Results}

The current quantitative outcomes are already sufficient to anchor the paper's central story. Averaged over the seven benchmarks, the proposed flow achieves 249.08\,MHz \fmax, compared with 231.91\,MHz for Vivado and 237.70\,MHz for \rsad. This corresponds to a 7.4\% \fmax gain over Vivado and a 4.8\% gain over \rsad. In terms of wirelength, the proposed flow reduces average \hpwl from 683{,}314 to 519{,}018 relative to Vivado, while remaining very close to \rsad (509{,}232 versus 519{,}018, only +1.9\%).

The timing metrics are consistent with the same story. Compared with Vivado, the proposed flow reduces the magnitude of average worst negative slack (\wns) from 0.71 to 0.39 and the magnitude of average total negative slack (\tns) from 1311.67 to 952.81. Compared with \rsad, it reduces the magnitude of average \wns from 0.61 to 0.39 and the magnitude of average \tns from 1245.90 to 952.81. These numbers suggest two points. First, the timing-driven objective is effective in moving the design toward better timing, which is consistent with the intended longest-wirelength control path. Second, the small \hpwl regression against \rsad indicates that the new method is not simply a wirelength minimizer; it is trading a very small amount of aggregate wirelength for stronger timing quality.

\begin{table}[t]
\caption{Average results over seven benchmarks.}
\label{tab:main}
\centering
\small
\begin{tabular}{lcccc}
\toprule
Method & HPWL & WNS & TNS & \fmax \\
\midrule
Vivado & 683{,}314 & -0.71 & -1311.67 & 231.91 \\
\rsad  & 509{,}232 & -0.61 & -1245.90 & 237.70 \\
AMFPlace2.0 + Vivado P\&R & 927{,}999 & -4.99 & -12837.49 & 117.46 \\
DSPlace + Vivado P\&R & 746{,}915 & -1.67 & -5634.44 & 191.72 \\
Ours   & 519{,}018 & -0.39 & -952.81 & 249.08 \\
\bottomrule
\end{tabular}
\end{table}

\begin{table*}[t]
\caption{Benchmark-wise \hpwl and \fmax comparison for all five methods.}
\label{tab:bench_hpwl_fmax}
\centering
\scriptsize
\resizebox{\textwidth}{!}{%
\begin{tabular}{l|rr|rr|rr|rr|rr}
\toprule
& \multicolumn{2}{c|}{Vivado} & \multicolumn{2}{c|}{\rsad} & \multicolumn{2}{c|}{AMFPlace 2.0 + Vivado P\&R} & \multicolumn{2}{c|}{DSPlace + Vivado P\&R} & \multicolumn{2}{c}{Ours} \\
Case & \hpwl & \fmax & \hpwl & \fmax & \hpwl & \fmax & \hpwl & \fmax & \hpwl & \fmax \\
\midrule
OS24x24\_W8 & 245232 & 263.71 & 177188 & 261.99 & 282005 & 136.91 & 292413 & 241.25 & 177441 & 262.54 \\
OS32x32\_W8 & 494792 & 239.06 & 356790 & 222.52 & 628927 & 121.86 & 543872 & 187.62 & 376203 & 244.92 \\
OS36x32\_W8 & 555410 & 202.51 & 568842 & 246.61 & 908625 & 116.55 & 650782 & 194.67 & 578173 & 247.71 \\
OS40x32\_W8 & 690842 & 230.41 & 448668 & 198.41 & 969454 & 124.35 & 781080 & 180.57 & 475704 & 244.56 \\
OS40x36\_W8 & 774885 & 238.55 & 540131 & 246.79 & 993061 & 116.89 & 880601 & 173.22 & 539512 & 250.50 \\
OS40x40\_W8 & 854961 & 233.26 & 597408 & 245.34 & 1704262 & 104.01 & 982803 & 184.26 & 593826 & 248.45 \\
OS44x40\_W8 & 1167074 & 215.84 & 875596 & 242.25 & 1009659 & 101.68 & 1096855 & 180.44 & 892266 & 244.86 \\
\midrule
AVG & 683314 & 231.91 & 509232 & 237.70 & 927999 & 117.46 & 746915 & 191.72 & 519018 & 249.08 \\
\bottomrule
\end{tabular}
\vspace{-4pt}
}
\end{table*}

\begin{table*}[t]
\caption{Benchmark-wise \wns and \tns comparison for all five methods.}
\label{tab:bench_wns_tns}
\centering
\scriptsize
\resizebox{\textwidth}{!}{%
\begin{tabular}{l|rr|rr|rr|rr|rr}
\toprule
& \multicolumn{2}{c|}{Vivado} & \multicolumn{2}{c|}{\rsad} & \multicolumn{2}{c|}{AMFPlace 2.0 + Vivado P\&R} & \multicolumn{2}{c|}{DSPlace + Vivado P\&R} & \multicolumn{2}{c}{Ours} \\
Case & \wns & \tns & \wns & \tns & \wns & \tns & \wns & \tns & \wns & \tns \\
\midrule
OS24x24\_W8 & -0.29 & -273.47 & -0.32 & -297.27 & -3.70 & -4584.19 & -0.55 & -485.23 & -0.31 & -351.68 \\
OS32x32\_W8 & -0.58 & -1298.78 & -0.89 & -2289.01 & -4.61 & -9662.71 & -1.73 & -4358.38 & -0.48 & -1455.69 \\
OS36x32\_W8 & -1.34 & -1537.04 & -0.46 & -1002.22 & -4.98 & -12272.79 & -1.54 & -4260.06 & -0.44 & -988.01 \\
OS40x32\_W8 & -0.74 & -1564.21 & -1.44 & -2285.96 & -4.44 & -12776.95 & -1.94 & -6675.55 & -0.49 & -1239.85 \\
OS40x36\_W8 & -0.59 & -1442.15 & -0.45 & -1115.90 & -4.96 & -13466.69 & -2.17 & -7940.15 & -0.39 & -1055.06 \\
OS40x40\_W8 & -0.39 & -39.24 & -0.18 & -22.46 & -6.01 & -20995.48 & -1.83 & -6782.90 & -0.13 & -12.35 \\
OS44x40\_W8 & -1.03 & -3026.82 & -0.53 & -1708.49 & -6.24 & -16103.59 & -1.94 & -8938.83 & -0.48 & -1567.06 \\
\midrule
AVG & -0.71 & -1311.67 & -0.61 & -1245.90 & -4.99 & -12837.49 & -1.67 & -5634.44 & -0.39 & -952.81 \\
\bottomrule
\end{tabular}
}
\end{table*}

\subsection{Interpretation}

The paper should interpret these results through the mechanism, not only through the percentages. Tables~\ref{tab:bench_hpwl_fmax} and~\ref{tab:bench_wns_tns} show that the proposed method is the most timing-effective approach among all five compared methods: it delivers the best average \fmax and the best average timing-slack metrics, while keeping \hpwl close to \rsad and substantially better than the three placement-and-route baselines built around generic flows. This is exactly the behavior expected from the method. Explicitly limiting the longest wirelength is a practical path to lowering \wns and improving \fmax. This is why \fmax should remain the first metric in the result discussion, while \hpwl serves as the supporting tradeoff metric. The role of \rsad should also be stated carefully: \rsad remains the stronger \hpwl-oriented predecessor because it exploits the OC structure effectively, and our method wins by adding timing awareness on top of that same structural foundation.

\subsection{Pending Insertions}

The next revision should add:
\begin{itemize}[leftmargin=*,topsep=2pt,itemsep=2pt]
\item one ablation on the timing-driven weight update,
\item one case-study figure showing a representative placement effect, and
\item a short paragraph discussing why \rsad can still be better in pure \hpwl on some cases while losing in \fmax.
\end{itemize}
